\subsection{Subcategory projection and exact twist support}
\label{sec:subcategory-projection}

Information about the center of a proper fusion subcategory imposes constraints
on individual restriction columns of the ambient center. It does not require
a braiding on the ambient fusion category, or an extension of a chosen
subcategory braiding.

\begin{proposition}\label{prop:central-subcategory-projection}
Let $\mathcal D\subset\mathcal C$ be a full fusion subcategory of a pivotal
fusion category, with the inherited pivotal structure. Let $P(V)$ be the sum
of all simple summands of $V$ belonging to $\mathcal D$.
If $(V,c)$ is an object of $Z(\mathcal C)$, then $P(V)$ has a canonical
$\mathcal D$-half-braiding. If $V$ is central simple with twist $\lambda$,
every simple constituent of $P(V)$ in $Z(\mathcal D)$ has twist $\lambda$.
Consequently, writing $F_{\mathcal D}:Z(\mathcal D)\to\mathcal D$ for the
forgetful functor, the $\mathcal D$-part of a central restriction column
belongs to the nonnegative integral span of the columns
$[F_{\mathcal D}(W)]$ with $\theta_W=\lambda$.
\end{proposition}
\begin{proof}
The complement of $\mathcal D$ is stable under tensoring on either side by
objects of $\mathcal D$. Indeed, if a simple $S\notin\mathcal D$ and
$X,T\in\mathcal D$ satisfy $\operatorname{Hom}(T,X\otimes S)\ne0$,
rigidity makes $S$ a constituent of $X^*\otimes T\in\mathcal D$, a
contradiction. Hence $P$ is a projection of $\mathcal D$-bimodule categories.
Put $W=P(V)$ and choose $i:W\to V$, $p:V\to W$ with $pi=1_W$.
For $X\in\mathcal D$ the maps
\[
 c^W_X=(1_X\otimes p)c_X(i\otimes1_X)
\]
are isomorphisms. Naturality and the half-braiding identity restrict to this
summand, so they make $W$ an object of $Z(\mathcal D)$.

It remains to check the twist, rather than merely the half-braiding.
Suppress associators and use the Drinfeld isomorphism
\[
 u_V=(\operatorname{ev}_V\otimes1_{V^{**}})
       (c_{V^*}\otimes1_{V^{**}})
       (1_V\otimes\operatorname{coev}_{V^*}):V\longrightarrow V^{**}.
\]
The duality identities give
\[
 \begin{split}
 (1_{V^*}\otimes p^{**})\operatorname{coev}_{V^*}
   &=(p^*\otimes1_{W^{**}})\operatorname{coev}_{W^*},\\
 \operatorname{ev}_V(p^*\otimes1_V)
   &=\operatorname{ev}_W(1_{W^*}\otimes p).
 \end{split}
\]
Using naturality of $c$ in its second argument between these two identities
gives $p^{**}u_Vi=u_W$. If $j$ denotes the pivotal isomorphism, then
$p^{**}j_Vi=j_W$. With the convention $\theta=u^{-1}j$, central simplicity
gives $u_V=\lambda^{-1}j_V$, and therefore
$u_W=\lambda^{-1}j_W$. Thus $\theta_W=\lambda1_W$.
Decomposing $W$ into central simples proves the integral-span assertion.
The opposite twist convention gives the same conclusion when used in both
centers.
\end{proof}

\begin{proposition}[Modular subcategory projection]
\label{prop:modular-subcategory-projection}
Suppose $\mathcal C$ is pseudo-unitary and $\mathcal D\subset\mathcal C$
is a full fusion subcategory, both with their positive spherical structures.
Let $\Pi:Z(\mathcal C)\to Z(\mathcal D)$ be the additive central
projection constructed above, and put
\[
 B_{WZ}=\dim\operatorname{Hom}_{Z(\mathcal D)}(W,\Pi Z).
\]
For the canonical normalized modular matrices of the two centers,
\begin{equation}\label{eq:subcategory-modular-intertwiner}
 S_{\mathcal D}B=BS_{\mathcal C},\qquad
 T_{\mathcal D}B=BT_{\mathcal C},\qquad
 B_{W,\mathbf1}=\delta_{W,\mathbf1}.
\end{equation}
In particular, each individual simple $W\in Z(\mathcal D)$ has the mass
and capacity bounds
\begin{equation}\label{eq:individual-central-projection-capacity}
 \sum_Z B_{WZ}d_Z
   =\frac{D_{\mathcal C}}{D_{\mathcal D}}d_W,
 \qquad
 B_{WZ}d_Z\le\frac{D_{\mathcal C}}{D_{\mathcal D}}d_W.
\end{equation}
The second inequality concerns one ambient central simple and one
subcategory central simple, even if several $W$ have the same twist and
forgetful image.
\end{proposition}
\begin{proof}
Regard $\mathcal M=\mathcal C$ as a
$\mathcal D$--$\mathcal C$ bimodule by tensor multiplication.
It is indecomposable, since its right $\mathcal C$ action alone generates
every simple from the unit. The full center of this module category is a
Lagrangian algebra
\[
 L\in Z(\mathcal D)\boxtimes Z(\mathcal C)^{\mathrm{rev}};
\]
see the module-category full-center construction
\cite[Section~6, Theorem~6.2]{DavydovCentre} and the correspondence
\cite[Proposition~4.8]{DMNO}. We identify its multiplicities explicitly.
For $W\in Z(\mathcal D)$ and $Z\in Z(\mathcal C)$, define bimodule
endofunctors on $\mathcal M$ by
\[
 L_W(X)=F_{\mathcal D}(W)\otimes X,
 \qquad R_Z(X)=X\otimes F_{\mathcal C}(Z).
\]
Their left and right module structures use the respective half-braidings.
The defining full-center adjunction identifies the coefficient of
$W\boxtimes Z^*$ with
$\dim\operatorname{Nat}_{\mathcal D-\mathcal C}(L_W,R_Z)$.
A right $\mathcal C$-module natural transformation is determined by its
component at the unit, a map
$f:F_{\mathcal D}(W)\to F_{\mathcal C}(Z)$.
Its left $\mathcal D$-linearity says precisely
\[
 c^Z_X(f\otimes1_X)=(1_X\otimes f)c^W_X
       \qquad(X\in\mathcal D).
\]
Because its source belongs to $\mathcal D$, $f$ factors uniquely through
$P(F_{\mathcal C}(Z))$. The displayed condition then says that this
factorization is a morphism $W\to\Pi Z$ in $Z(\mathcal D)$.
Consequently
\[
 L\cong\bigoplus_{W,Z}B_{WZ}\,W\boxtimes Z^*.
\]
This fixes the dual convention: when $\mathcal D=\mathcal C$, $B$ is the
identity and $L$ is the usual diagonal Lagrangian algebra.

The class of a Lagrangian algebra is fixed by the canonical normalized
$S$ and $T$. Indeed, the Morita description identifies it with induction
of the unit in a center, where this is
\cite[Theorem~2.5, Proposition~2.8]{Ostrik3}; the positive spherical
structures make the center equivalence ribbon. Let $C_{\mathcal C}$ be
the duality permutation of the ambient center. The coefficient matrix
of $L$ is $BC_{\mathcal C}$, and the reversed center has matrix
$\overline{S_{\mathcal C}}=S_{\mathcal C}^{-1}$ and inverse twists.
Thus invariance is
\[
 S_{\mathcal D}BC_{\mathcal C}S_{\mathcal C}^{-1}
       =BC_{\mathcal C},\qquad
 T_{\mathcal D}BC_{\mathcal C}T_{\mathcal C}^{-1}
       =BC_{\mathcal C}.
\]
Duality commutes with both matrices, so these equalities give the first
two identities in \eqref{eq:subcategory-modular-intertwiner}.
Projection sends the ambient central unit to the subcategory central
unit, proving the last identity. Finally compare the ambient-unit column
of $S_{\mathcal D}B=BS_{\mathcal C}$ and use
$S_{\mathcal D,W0}=d_W/D_{\mathcal D}$ and
$S_{\mathcal C,Z0}=d_Z/D_{\mathcal C}$.
Positivity of all summands proves the individual capacity inequality.
\end{proof}

\begin{corollary}[Disjoint central projection rows]
\label{cor:disjoint-central-projection}
In the setting of Proposition~\ref{prop:modular-subcategory-projection},
suppose every ambient central simple projects either to zero or to a
single simple of $Z(\mathcal D)$ with multiplicity one. Then, for a
positive integer $r$,
\[
 BB^{\mathsf T}=rI.
\]
Thus $B^{\mathsf T}/\sqrt r$ is an isometric intertwining embedding of
the entire canonical modular representation of $Z(\mathcal D)$ into
that of $Z(\mathcal C)$. Its $T$-spectrum can be subtracted before the
finite-quotient representation test. Any signed charge-conjugation
bound for the complement is obtained by subtracting this known
subrepresentation's signed trace from the corresponding full-space
bound; it need not be nonnegative.
\end{corollary}
\begin{proof}
The hypothesis says each column of $B$ is zero or a standard basis
vector. Hence $BB^{\mathsf T}$ is diagonal, with entries equal to the
numbers of simples in the projection fibers. Each entry is positive
by \eqref{eq:individual-central-projection-capacity}. The normalized
modular representations are unitary, so the intertwining identities
imply that $BB^{\mathsf T}$ commutes with $S_{\mathcal D}$ and
$T_{\mathcal D}$. Comparing its $W,\mathbf1$ entries and using
$S_{\mathcal D,W\mathbf1}=d_W/D_{\mathcal D}>0$ shows that all diagonal
entries equal the same positive integer $r$. The embedding and the
subtraction statements follow by taking adjoints. The orthogonal
complement is a representation space, but is not required to have a
simple basis permuted by charge conjugation.
\end{proof}

The disjoint-column hypothesis must be verified. The exact categorical
control \filename{research/round5/unification/subcategory_intertwiner_control.py}
derives $S$, $T$ and $B$ from finite group character sums for
$\mathrm{Vec}(C_2),\mathrm{Vec}(C_3)\subset\mathrm{Vec}(S_3)$.
Both rectangular modular identities and all individual masses hold.
The second projection has rank five although its target center has rank
nine, illustrating why a general subcategory inclusion alone does not
justify subtracting the entire target representation. Mutating one
restriction multiplicity violates the checked $S$-identity.

For an Ising subcategory, the two twist-minus-one simples of its center
are $\mathbf1\boxtimes g$ and $g\boxtimes\mathbf1$, each of dimension
one. Hence every high ambient central simple whose restriction contains
$g$ has dimension at most $D_{\mathcal C}/4$. This is stronger than the
combined twist-minus-one mass $D_{\mathcal C}/2$.
For $\mathcal D=\mathrm{Vec}$, the same proposition says that the
multiplicity vector of $I_{\mathcal C}(\mathbf1)$ is fixed by $S$ and $T$;
thus the full modular representation contains a trivial constituent.

The projection also gives exact, twist-resolved dimension identities.
Here and below spherical structures are inherited, so the centers carry their
standard ribbon structures.

\begin{corollary}\label{cor:subcategory-twist-masses}
Let $\mathcal C$ be spherical and $\mathcal D\subset\mathcal C$ a full fusion
subcategory. Put $D_{\mathcal C}=\dim(\mathcal C)$ and
$D_{\mathcal D}=\dim(\mathcal D)$. For a simple $X\in\mathcal D$ and a root
of unity $\lambda$,
\[
 \sum_{\substack{Z\in\operatorname{Irr}Z(\mathcal C)\\\theta_Z=\lambda}}
       \dim(Z)[F_{\mathcal C}(Z):X]
 =\frac{D_{\mathcal C}}{D_{\mathcal D}}
   \sum_{\substack{W\in\operatorname{Irr}Z(\mathcal D)\\\theta_W=\lambda}}
       \dim(W)[F_{\mathcal D}(W):X].
\]
An empty right-hand side is zero. In the pseudo-unitary case all weights are
positive and dimensions can be replaced by Frobenius--Perron dimensions.
\end{corollary}
\begin{proof}
Choose $N$ divisible by the orders of all twists in $Z(\mathcal D)$.
Proposition~\ref{prop:central-subcategory-projection} shows that every central
simple contributing to the left-hand side has twist in $\mu_N$.
The indicator formula \cite{NS} gives, for $1\le n<N$,
\[
 D_{\mathcal C}\nu_n(X)
 =\sum_Z\theta_Z^n\dim(Z)[F_{\mathcal C}(Z):X].
\]
The same formula holds in $\mathcal D$, and the indicators agree: their
defining pivotal rotations act on the same spaces
$\operatorname{Hom}(1,X^{\otimes n})$. For $n=0$, the corresponding equality
is $\dim I_{\mathcal C}(X)=D_{\mathcal C}\dim(X)$, and similarly in
$\mathcal D$. Fourier inversion on $\mu_N$ proves the assertion.
\end{proof}

\begin{corollary}\label{cor:fibonacci-ising-projection}
In the pseudo-unitary setting the following subcategory constraints hold.
\begin{enumerate}
\item For a Fibonacci subcategory with basis $(1,a)$, $a^2=1+a$, write
$\varphi=\operatorname{FPdim}(a)=(1+\sqrt5)/2$. A nonzero projected
column $(r,s)$ has either twist $1$ and $r\ge s$, or $r=0$ and twist
$\zeta_5^2$ or $\zeta_5^{-2}$. For each of the latter twists, the
$a$-mass in Corollary~\ref{cor:subcategory-twist-masses} is
$D_{\mathcal C}\varphi/(\varphi+2)$; the twist-one $a$-mass is
$D_{\mathcal C}\varphi^2/(\varphi+2)$.
\item For an Ising subcategory with basis $(1,g,\sigma)$, a nonzero projected
column $(r,s,t)$ has one of the following forms: twist $1$ with $t=0$ and
$r\ge s$; twist $-1$ with $r=t=0$; or one of four Ising twists
$\{z,-z,z^{-1},-z^{-1}\}$ with $r=s=0$, where $z$ is an appropriate
primitive sixteenth root. The $g$-masses at $1$ and $-1$ are each
$D_{\mathcal C}/2$. The $\sigma$-mass at each of its four twists is
$D_{\mathcal C}\sqrt2/4$.
\end{enumerate}
In particular, a high column, whose unit coordinate is zero, cannot combine
$g$ and $\sigma$ in its Ising projection, or contain either $a$, $g$ or
$\sigma$ at twist one in the respective cases.
\end{corollary}
\begin{proof}
The rank-two and Ising classifications \cite{Ostrik2,EGNO} provide
nondegenerate braidings on these subcategories. For either such category
$\mathcal D$, its center is
$\mathcal D\boxtimes\mathcal D^{\mathrm{rev}}$ and the restriction of
$(U,V)$ is $U\otimes V$.
For Fibonacci the four columns are
\[
 (1,0),\quad(0,1),\quad(0,1),\quad(1,1),
\]
with twists $1,z,z^{-1},1$. Positivity gives the usual Fibonacci modular
$S$-matrix; balancing gives $z^2+\varphi z+1=0$, hence
$z=\zeta_5^{\pm2}$. The center dimensions are respectively
$1,\varphi,\varphi,\varphi^2$.
For Ising, the twist-one columns are $(1,0,0)$ twice and $(1,1,0)$ once;
the twist-$-1$ columns are $(0,1,0)$ twice; and each of
$z,-z,z^{-1},-z^{-1}$ has the column $(0,0,1)$ once. Their dimensions are
$1,2,\sqrt2$ as determined by these columns. Taking their additive spans
and applying the preceding corollary gives every assertion.
All inherited Ising spherical choices must be retained. If the second
indicator of $\sigma$ is fixed, the identity
$z^2+z^{-2}=\nu_2(\sigma)\sqrt2$ selects the appropriate set of four twists.
\end{proof}

\paragraph{Computational use.}
The criteria first restrict the allowed high-column/twist pairs and the
twist-one low columns. They then add the exact mass identities to the
integer linear equations. When using normalized relative twist traces,
one projects the identities to the chosen field, combining coincident
trace types; one must not confuse a projected trace with an actual twist.
These are necessary conditions only. General exclusions still require the
separately checked general-to-pseudo-unitary reduction.
